\section{Conclusion}

The Bayesian Detection Score (BDS) converts a redistricting ensemble percentile
into a court-interpretable posterior probability. The key advantages over the
standard frequentist percentile:

\begin{enumerate}
  \item \textbf{Uncertainty quantification}: BDS incorporates the ensemble's
    ESS rather than treating $N$ nominal plans as $N$ independent observations.
  \item \textbf{Probabilistic interpretation}: ``BDS = 0.97'' is directly
    interpretable as ``97\% probability the plan is in the extreme 5\%,'' which
    is what courts want to know.
  \item \textbf{Geography absorption}: the ensemble prior absorbs geographic
    effects, so the posterior measures only the residual above geographic constraints.
  \item \textbf{Threshold clarity}: BDS $> 0.95$ as the evidentiary threshold
    has a clear probabilistic meaning, analogous to $p < 0.05$ but more honest
    about the uncertainty.
\end{enumerate}

Applied to five study states, the BDS confirms strong Bayesian evidence that
NC, WI, GA, and PA enacted plans are genuinely extreme (BDS $> 0.97$), while
the \textsc{bisect} algorithmic plans are consistent with neutral random draws
(BDS $< 0.10$). Texas presents an intermediate case (BDS = 0.89), reflecting
its geographic complexity.

The BDS is implemented as a three-line computation from standard statistical
libraries:
\begin{verbatim}
from scipy.stats import beta
bds = beta.cdf(0.05, p_hat*ess + 1, (1 - p_hat)*ess + 1)
\end{verbatim}
No new redistricting-specific software is required; the computation runs on
any GerryChain ensemble output with an ESS estimate from G.4.
